% set default size and document class
\documentclass[12pt]{report}

\usetheme{Berkeley}
\usecolortheme{seahorse}
\usebackgroundtemplate{\includegraphics[width=\paperwidth]{../../paper}}

\beamertemplatenavigationsymbolsempty

\addtobeamertemplate{navigation symbols}{}{
\color{black}\normalsize\textbf\insertframenumber}


\singlespacing

\newcommand{\hrq}[0]{\rule{\textwidth/4}{1pt}}

\begin{titlepage}
\cnam
\begin{center}
{\bfseries

\begin{large}

CONSERVATOIRE NATIONAL DES ARTS ET MÉTIERS

CENTRE RÉGIONAL ASSOCIÉ DE\\
NOUVELLE-AQUITAINE

\hrq

\goal

\vfill
par
\vfill

Marc BENINCA

\hrq

\begin{Large}\fulltitle\end{Large}

Soutenu le 12 Mai 2020

\hrq

JURY

\end{large}
\vspace{1em}
\begin{small}

\begin{tabular}[t]{rllll}
PRÉSIDENT : & Mme & Catherine & BRULATOUT & \textit{\begin{tabular}[t]{@{}l@{}}Maître de conférences associé\\Université de Bordeaux\end{tabular}} \\
MEMBRES : & Mme & Elisabeth & ABRIVAT & \textit{\begin{tabular}[t]{@{}l@{}}Maître de conférences associé\\Université de Bordeaux\end{tabular}} \\
%& ? & Prénom & NOM & \textit{\begin{tabular}[t]{@{}l@{}}Fonction\\Organisme\end{tabular}}
\end{tabular}

\end{small}

}
\end{center}
\end{titlepage}


\onehalfspacing

\ml
{\def\thanks{Acknowledgements}}
{\def\thanks{Remerciements}}
\chapter*{\thanks}
\addcontentsline{toc}{chapter}{\thanks}

\pagebreak

\printglossary[type=\acronymtype,
title=\ml{Acronyms}{Abréviations}]

\pagebreak

\ml
{\def\gloss{Glossary}}
{\def\gloss{Glossaire}}
\printglossary
[title=\gloss]

\newglossaryentry{exp}{
name={expression},
plural={expressions},
description={\ml
{english description}
{description en français}},
}

\pagebreak

\cnam

\begin{large}\fulltitle\end{large}

\renewcommand{\contentsname}{Plan}

\cftsetindents{section}{1em}{1.5em}
\cftsetindents{subsection}{2.5em}{2.5em}
\cftsetindents{subsubsection}{5em}{3.25em}
\renewcommand{\cftsecleader}{\cftdotfill{\cftdotsep}}
%\renewcommand{\cftsubsecleader}{\hfill}

\tableofcontents

\pagebreak

\section{Introduction}

Les \gls{tic} ont fait émerger des \gls{si} de grande envergure, qui génèrent
continuellement une quantité globale toujours plus importante de données,
communément appelée \gls{bd}.

Les volumes démesurés de données manipulés par ces \gls{si} sont clairement
incompatibles avec les principes classiques des \gls{sgbd}.
De nouvelles méthodologies ont donc dû être créées afin de pouvoir
analyser puis exploiter ces \gls{bd}.

\subsection{Machine learning}

Le \gls{ml} regroupe les différentes nouvelles approches méthodologiques
permettant de faire ressortir une compréhension des \gls{bd}, en extraire
des \glspl{ds} exploitables puis, à partir de cet apprentissage,
être en mesure de comprendre de nouvelles données.

Les 2 défis à relever pour être pertinents dans cette démarche :
\begin{itmz}
\item{ne pas surentraîner (spécialiser) l’apprentisseur automatique\\
pour conserver une capacité de généralisation}
\item{pour chaque type de problèmes à résoudre,\\
choisir les algorithmes les plus pertinents}
\end{itmz}

\subsection{Méthodes}

\subsubsection{Méthodes non supervisées}

Ces premières méthodes permettent de travailler avec un ensemble d’éléments
ne disposant pas d’étiquetage préalable, elles mettent ainsi en place
des mécanismes permettant de faire du \gls{clustering} de ces éléments.

Exemples d’algorithmes :
\begin{itmz}
\item{k-means (moyennes)}
\item{réduction de dimensionnalité}
\item{réseaux de neurones}
\item{analyse des composants principaux ou indépendants}
\item{modèles de distribution}
\item{classification hiérarchique}
\item{\gls{clustering} par décalage moyen}
\item{Apriori}
\end{itmz}

\subsubsection{Méthodes semi-supervisées}

Ces méthodes autorisent l’utilisation d’un ensemble d’éléments hétérogène,
certains étant ayant déjà été étiquetés au préalable, alors que d’autres non.
Un mélange de différentes techniques s’avère donc être nécessaire.

Exemples d’algorithmes :
\begin{itmz}
\item{\gls{hbos}}
\item{forêts isolées}
\item{autoencodeurs}
\end{itmz}

\subsubsection{Méthodes supervisées}

Ces méthodes utilisent 2 sous-ensembles d’éléments intégralement étiquetés,
ayant une valeur de sortie pour plusieurs valeurs de variables en entrée.
Les données d’apprentissage permettent un entraînement préalable du modèle.
Les données de test permettent ensuite d’en évaluer la pertinence.

Le modèle permet enfin, à partir de nouvelles valeurs en entrée,
de fournir des valeurs de sortie prédictives de 2 types :
\begin{itmz}
\item{régression → la valeur est un nombre}
\item{classification → la valeur est une catégorie, parmi un ensemble défini}
\end{itmz}

Exemples d’algorithmes :
\begin{itmz}
\item{régressions linéaire, logistique ou vectorielle}
\item{arbres de régression, classification}
\item{k-\gls{nn} (plus proches voisins)}
\item{classificateur Naïve Bayes}
\item{réseaux de neurones}
\item{\gls{svm}}
\end{itmz}

Ce sont ces \gls{svm} qui feront l’objet d’étude du présent document.
En seront passés en revue les différents principes techniques sous-jacents,
un exemple d’application, ainsi qu’un regard critique sur ce type d’algorithmes.

\pagebreak


\ml
{\chapter{Problem: maintenance of operating systems}}
{\chapter{Problématique : maintenance des systèmes d’exploitation}}

\begin{itmz}
\item{\ml{updates}
{mises à jour}}
\item{\ml{unavailability}
{indisponibilité}}
\item{\ml{backups policy}
{politique de sauvegardes}}
\item{\ml{testing backups}
{tests des sauvegardes}}
\item{\ml{snapshots}
{instantanés}}
\item{\ml{restorations}
{restaurations}}
\item{\ml{configuration recipes}
{recettes de configuration}}
\end{itmz}

\ml
{\section{File systems, installed on partitions, with write access}}
{\section{Systèmes de fichiers, installés sur partitions, avec accès en écriture}}

\ml
{\subsection{Conventional file systems}}
{\subsection{Système de fichiers conventionnel}}

\begin{itmz}
\item{ext2}
\item{ext3}
\item{ext4}
\item{jfs}
\item{xfs}
\end{itmz}

\ml{Pros}{Avantages}:
\begin{itmz}
\item{\ml{\todo}
{\todo}}
\end{itmz}

\ml{Cons}{Inconvénients}:
\begin{itmz}
\item{\ml{\todo}
{\todo}}
\end{itmz}

\ml
{\subsection{File systems managed with configuration recipes}}
{\subsection{Système de fichiers géré par des recettes configuration}}

\begin{itmz}
\item{ansible}
\item{chef}
\item{puppet}
\end{itmz}

\ml{Pros}{Avantages}:
\begin{itmz}
\item{\ml{\todo}
{\todo}}
\end{itmz}

\ml{Cons}{Inconvénients}:
\begin{itmz}
\item{\ml{\todo}
{\todo}}
\end{itmz}

\ml
{\subsection{File systems supporting snapshots}}
{\subsection{Système de fichiers avec gestion d’instantanés}}

\begin{itmz}
\item{btrfs}
\item{zfs}
\end{itmz}

\ml{Pros}{Avantages}:
\begin{itmz}
\item{\ml{\todo}
{\todo}}
\end{itmz}

\ml{Cons}{Inconvénients}:
\begin{itmz}
\item{\ml{\todo}
{\todo}}
\end{itmz}

\ml
{\section{Live images, installationless, with read access}}
{\section{Images autonomes, sans installation, avec accès en lecture seule}}

\ml
{\subsection{Boot without persistence}}
{\subsection{Amorçage sans gestion de persistance}}

\ml{Pros}{Avantages}:
\begin{itmz}
\item{\ml{\todo}
{\todo}}
\end{itmz}

\ml{Cons}{Inconvénients}:
\begin{itmz}
\item{\ml{\todo}
{\todo}}
\end{itmz}

\ml
{\subsection{Boot with persistence}}
{\subsection{Amorçage avec gestion de persistance}}

\ml{Pros}{Avantages}:
\begin{itmz}
\item{\ml{\todo}
{\todo}}
\end{itmz}

\ml{Cons}{Inconvénients}:
\begin{itmz}
\item{\ml{\todo}
{\todo}}
\end{itmz}

\ml
{\section{Existing alternatives}}
{\section{Alternatives existantes}}

\ml
{\chapter{Proposal: an incremental live workflow}}
{\chapter{Proposition : un fonctionnement autonome incrémental}}

\ml{Pros}{Avantages}:
\begin{itmz}
\item{\ml{reboot = restore}
{redémarrage = restauration}}
\item{\ml{update = backup}
{mise à jour = sauvegarde}}
\item{\ml{separation of system and data}
{séparation système et données}}
\end{itmz}

\ml{Cons}{Inconvénients}:
\begin{itmz}
\item{\ml{exhaustive manual procedure}
{maintenance manuelle exhaustive}}
\end{itmz}

\ml
{\section{Implement the workflow}}
{\section{Mettre en œuvre le fonctionnement}}

\ml
{\subsection{Mirrors of official and third-party repositories}}
{\subsection{Miroirs de dépôts officiels et éditeurs tiers}}

\ml{Pros}{Avantages}:
\begin{itmz}
\item{\ml{\todo}
{\todo}}
\end{itmz}

\ml{Cons}{Inconvénients}:
\begin{itmz}
\item{\ml{\todo}
{\todo}}
\end{itmz}

\ml
{\subsubsection{Local synchronization}}
{\subsubsection{Synchronisation locale}}

\paragraph{apt-mirror}

\begin{itmz}
\item{\ml{translations}
{traductions} (Translation-*.bz2)}
\item{\ml{architecture independant contents}
{contenus indépendants de l’architecture} (Contents-all.gz)}
\item{InRelease \ml{with some third-party repositories}
{avec certains dépôts éditeurs}}
\end{itmz}

\paragraph{debmirror}

\paragraph{ftpsync}

\ml
{\subsubsection{Select useful architectures}}
{\subsubsection{Choisir les architectures utiles}}

\paragraph{amd64}

\paragraph{arm64}

\paragraph{armhf}

\paragraph{i386}

\ml
{\subsubsection{Check integrity}}
{\subsubsection{Vérifier l’intégrité}}

\ml{Pros}{Avantages}:
\begin{itmz}
\item{\ml{avoid errors during future package installations}
{éviter des erreurs lors de futures installations de paquets}}
\end{itmz}

\ml{Cons}{Inconvénients}:
\begin{itmz}
\item{\ml{no tool exists}
{aucun outil n’existe}}
\end{itmz}

\ml
{\subsection{Building a live file system}}
{\subsection{Construction d’un système de fichiers autonome}}

\paragraph{Debian GNU/Linux}

\ml
{\subsubsection{Install specific packages}}
{\subsubsection{Installer les paquets spécifiques}}

\paragraph{\ml{Bare metal}{Machine physique}}

\paragraph{\ml{Virtual machine}{Machine virtuelle}}

\paragraph{\ml{Container}{Conteneur}}

\ml
{\subsubsection{Create a minimal file system base}}
{\subsubsection{Créer un système de fichiers minimal}}

\paragraph{debootstrap}

\ml
{\subsubsection{Turn a system into a systems factory}}
{\subsubsection{Équiper un système pour en fabriquer d’autres}}

\ml
{\subsubsection{Turn a file system into a live one}}
{\subsubsection{Rendre un système de fichiers autonome}}

\paragraph{live-boot}

\paragraph{update-initramfs}

\ml
{\subsubsection{Install additional packages}}
{\subsubsection{Installer des paquets supplémentaires}}

\ml
{\subsubsection{Data persistence through symbolic links}}
{\subsubsection{Déport des données persistantes avec liens symboliques}}

\ml
{\subsection{Encapsulate in an image file}}
{\subsection{Encapsuler dans un fichier image}}

\ml
{\subsubsection{Use a format suited for read-only mounting}}
{\subsubsection{Utiliser un format adapté au montage en lecture}}

\paragraph{SquashFS}

\ml
{\subsubsection{Choose a compression algorithm}}
{\subsubsection{Choisir un algorithme de compression}}

\paragraph{gzip}

\paragraph{lzma}

\paragraph{lzo}

\paragraph{lz4}

\paragraph{xz}

\paragraph{zstd}

\ml
{\subsubsection{Encapsulate in a hybrid image file}}
{\subsubsection{Encapsuler dans un fichier image hybride}}

\paragraph{ISO}

\ml
{\subsection{Secure the image file produced}}
{\subsection{Sécuriser un fichier image produit}}

\ml
{\subsubsection{Compute an integrity checksum}}
{\subsubsection{Calculer une somme de contrôle d’intégrité}}

\paragraph{SHA-256}

\paragraph{SHA-512}

\ml
{\subsubsection{Sign to certify authenticity}}
{\subsubsection{Signer pour certifier l’authenticité}}

\ml
{\subsection{Boot secure image files}}
{\subsection{Amorcer des fichiers images sécurisés}}

\ml
{\subsubsection{Create standalone boot images}}
{\subsubsection{Créer des images de démarrage}}

\paragraph{GRUB}

\paragraph{BIOS}

\paragraph{UEFI}

\paragraph{Secure boot}

\ml
{\subsubsection{Create a boot menu}}
{\subsubsection{Créer un menu de démarrage}}

\paragraph{grub.cfg}

\paragraph{loopback}

\paragraph{squash4}

\paragraph{iso9660}

\ml
{\subsubsection{Check integrity and authenticity}}
{\subsubsection{Vérifier intégrité et authentiticté}}

\paragraph{gcry\_sha256}

\paragraph{gcry\_sha512}

\paragraph{pgp}

\ml
{\subsubsection{Load images in random access memory}}
{\subsubsection{Charger des images en mémoire vive}}

\paragraph{overlayfs}

\ml
{\subsection{Incremental updating}}
{\subsection{Mise à niveau incrémentale}}

\ml
{\subsubsection{Create a new image file}}
{\subsubsection{Fabriquer une nouvelle image}}

\ml
{\subsubsection{Avoid an unnecessary reboot}}
{\subsubsection{Éviter un redémarrage facultatif}}

\ml{Pros}{Avantages}:
\begin{itmz}
\item{\ml{no down time}
{aucune indisponibilité}}
\item{\ml{just replay the modifications on the system in memory}
{répéter simplement les modifications sur le système en mémoire}}
\end{itmz}

\ml{Cons}{Inconvénients}:
\begin{itmz}
\item{\ml{\todo}
{\todo}}
\end{itmz}

\ml
{\subsubsection{Reduce the duration of a mandatory reboot}}
{\subsubsection{Réduire la durée d’un redémarrage obligatoire}}

\paragraph{kexec-tools}

\ml
{\section{Automate the workflow}}
{\section{Automatiser le fonctionnement}}

\ml
{\subsection{Check integrity of local repositories}}
{\subsection{Vérifier l’intégrité des dépôts locaux}}

\ml
{\subsection{Build complete live file systems from scratch}}
{\subsection{Construire des systèmes de fichiers autonomes complets}}

\ml
{\subsection{Create new files by updating existing images}}
{\subsection{Créer de nouveaux fichiers par mise à jour d’images}}

\ml
{\subsection{Generate a boot menu on-the-fly}}
{\subsection{Générer un menu de démarrage à la volée}}

\ml
{\chapter{Results}}
{\chapter{Résultats}}

\section{Conclusion}

Nombreux outils logiciels implémentant les \gls{svm},
(libSVM, SVM-Light, SVMTorch) disponibles pour différents langages.

Nativement bons linéairement, noyau permet apprentissage non-linéaire.

Noyau élément clé, \gls{kf} pas facile à choisir.

Il faut savoir rester patient pour les temps d’entraînement.

\pagebreak


\appendix

\ml
{\chapter{Appendix}}
{\chapter{Annexe}}
\pagebreak

\pagebreak

\def\figures{Figures}
\renewcommand{\listfigurename}{\figures}
\listoffigures
\addcontentsline{toc}{chapter}{\figures}

\def\tables{\ml{Tables}{Tableaux}}
\renewcommand{\listtablename}{\tables}
\listoftables
\addcontentsline{toc}{chapter}{\tables}


\singlespacing

\chapter*{}
\addcontentsline{toc}{chapter}{\ml{Summaries}{Résumés}}

\title,\\
\subtitle.

\ml{CNAM Engineer thesis}{Mémoire d’Ingénieur CNAM}, Bordeaux 2021.

\summaries

\thispagestyle{empty}


\end{document}
